\documentclass[a4paper, 12pt]{article}

\usepackage[utf8]{inputenc}
\usepackage[T1]{fontenc}
\usepackage{textcomp}
\usepackage{amssymb}
\usepackage{newtxtext} \usepackage{newtxmath}
\usepackage{amsmath, amssymb}
\newtheorem{problem}{Problem}
\newtheorem{example}{Example}
\newtheorem{lemma}{Lemma}
\newtheorem{theorem}{Theorem}
\newtheorem{problem}{Problem}
\newtheorem{example}{Example} \newtheorem{definition}{Definition}
\newtheorem{lemma}{Lemma}
\newtheorem{theorem}{Theorem}
\DeclareMathAlphabet{\mathcal}{OMS}{cmsy}{m}{n}

\begin{document}

\section*{Mathematical background}

The Hilbert transform of a function $f$ is defined as 

$$\mathcal{H}\left\{ f(x) \right\}  = \frac{1}{\pi} \text{P.V.} \int_{-\infty}^{\infty} \frac{f(t)}{x - t} dt,$$

where P.V. denotes the Cauchy principal value of the integral and serves the
purpose of ensuring that the integral converges. Clearly, the transform is
nothing but a convolution of $f$ with $1 / \pi t$. The transform effectively 
shifts the phase of each frequency component of the function by $\pi / 2$
radians, i.e. it satisfies that  

\begin{equation*}
    \mathcal{F} \{ \mathcal{H}\{ f(x) \} \} = -i \cdot \text{sgn}(\omega) \cdot
    \mathcal{F}\{ f(x) \}
\end{equation*}

This means that the function

\begin{equation*}
    z(t) = f(t) + i \cdot \mathcal{H}\{ f(t) \}
\end{equation*}

has no negative frequency components (they cancel out). The function $z(t)$ is
called the analytic signal of $f(t)$, and the magnitude of $z(t)$ is called the
envelope of $f(t)$. The analytic signal allows us to extract the instantaneous
amplitude and phase of a signal: 

\begin{equation*}
    A(t) = |z(t)| = \sqrt{f(t)^2 + \mathcal{H}\{ f(t) \}^2}
\end{equation*}

\begin{equation*}
    \phi(t) = \arg(z(t)) = \arctan\left( \frac{\mathcal{H}\{ f(t) \}}{f(t)}
    \right)
\end{equation*}

From the instantaneous phase we derive the instantaneous frequency as the time
derivative of the phase:

\begin{equation*}
    \omega(t) = \frac{d\phi(t)}{dt}
\end{equation*}

The capacity to provide the \textit{instantaneous} frequency and amplitude is
what distinguishes the Hilbert transform from the Fourier transform, which only
provides the average frequency content of a signal.

\subsection*{Empirical mode decomposition}

Empirical mode decomposition (EMD) is a method for decomposing a signal into a
set of intrinsic mode functions (IMFs) that represent the different oscillatory
modes present in the signal. The EMD process involves the following steps:

1. Identify the local maxima and minima of the signal.
2. Interpolate the local maxima to create an upper envelope and the local minima
to create a lower envelope.
3. Compute the mean of the upper and lower envelopes to obtain a mean envelope. 
4. Subtract the mean envelope from the original signal to obtain a candidate
IMF.
5. Check if the candidate IMF satisfies the following conditions:
   a. The number of extrema and the number of zero crossings must either be equal
   or differ at most by one.
   b. The mean value of the envelope defined by the local maxima and the
   envelope defined by the local minima must be zero at any point.
6. If the candidate IMF does not satisfy the conditions, repeat steps 1-5 on
the candidate IMF until an IMF is obtained.
7. Subtract the obtained IMF from the original signal to get a residual signal.
8. Repeat the entire process on the residual signal until it becomes a monotonic
function or contains only one extremum.











\end{document}



